% =======================================================
\section{Related Work}
\label{sec:related}
% =======================================================

\paragraph{Variational regularization.} Tikhonov regularization and total variation (TV) \citep{rudin1992tv}
remain standard baselines for ill-posed inverse problems because they are convex, well understood and
simple to tune. TV is a strong prior for piecewise-constant images, which makes it a competitive baseline
for digit-like data and a natural reference for CT.

\paragraph{Learned CT reconstruction.} Deep networks have been used to post-process filtered
back-projection images \citep{jin2017}, to unroll iterative reconstruction \citep{adler2018}, and, more
recently, as generative priors for sparse-view CT and accelerated MRI \citep{song2022}. On the LoDoPaB-CT
benchmark, FBP followed by a U-Net improves the test PSNR from $30.5$ to $35.8$\,dB over FBP
\citep{leuschner2021}. These methods are trained for a given operator and noise level. We instead study the
Plug-and-Play family, in which the prior is trained once and the data-fidelity step stays explicit, so that
the prior can be compared with classical regularizers inside the same iteration; our conclusions concern
this family and a small denoiser, not end-to-end networks.

\paragraph{Plug-and-Play and denoiser-based priors.}
\citet{venkatakrishnan2013plugandplay} proposed replacing the proximal operator of a variational regularizer by
an off-the-shelf denoiser inside an ADMM or proximal-gradient iteration; regularization by denoising
\citep{romano2017} is a related formulation, and \citet{kamilov2023} review the field. CNN denoisers were
introduced as plug-in priors by \citet{zhang2017ircnn}. Convergence results require assumptions on the denoiser,
such as boundedness or a Lipschitz constraint \citep{chan2016pnp,ryu2019pnp}; an unconstrained CNN denoiser
comes with no such guarantee, which is relevant to the peak-then-decay behaviour we observe in
Section~\ref{sec:experiments}.

\paragraph{Noise-conditioned denoisers.} DnCNN \citep{zhang2017dncnn} uses residual learning for Gaussian
denoising and, in its blind form, is trained over a range of noise levels. FFDNet \citep{zhang2018ffdnet}
instead receives the noise level as an input. We use a DnCNN-style residual network with the noise level
as an additional input channel, so that one network serves every denoising strength $\sigma_d$ tried in the
PnP iteration, where $\sigma_d$ is a tuned hyperparameter.

\paragraph{Instabilities, robustness and distribution shift.} \citet{antun2020instabilities} showed, for
MRI and CT networks, that tiny perturbations can cause severe artefacts, that small structural changes
may be missed, and that sparsity-based benchmarks including TV were unaffected in their examples.
\citet{darestani2021shift} measured the robustness of trained and untrained reconstruction methods in
accelerated MRI and found that both suffer similarly from distribution shifts and that the better-performing
deep methods do not lose robustness. Our shifts are natural (another image class, a different kind of
phantom) rather than adversarial, our prior is a small PnP denoiser, and the shifts are larger than within-domain
variation; we report a reversal of the advantage over TV for a mismatched prior (Section~\ref{sec:discussion}),
and we do not claim that it holds for other learned methods.

\paragraph{Hallucinations, the inverse crime and task-based assessment.}
\citet{bhadra2021} define hallucinations in tomographic reconstruction as errors in image components that the
measurements cannot stably determine, which can come only from the prior. In a stylised 2D MRI setting
they found that a U-Net produced more null-space hallucinations on out-of-distribution data, where its advantage
over non-learned methods disappeared. We adopt this view and extend the decomposition to the weakly determined
subspace of a finite-precision operator (Section~\ref{sec:subspace}). Our setting differs in the modality (CT),
the prior (a PnP denoiser), the controls (a prior trained on the shifted distribution, denoiser seeds, a TV
reference with the same iteration budget) and the evidence (a lesion-insertion test, and real data over a grid of
doses and angle counts). Generating data with the discrete model that is later inverted is known as the inverse
crime \citep{wirgin2004,kaipio2005} and tends to overstate the accuracy of methods that exploit the exact model;
we remove it in Section~\ref{sec:crimefree}. Finally, we follow the task-based view of image quality
\citep{barrett1993,barrett2004}: fidelity metrics such as PSNR need not predict whether a feature of interest
survives reconstruction, so we also measure the detection of inserted lesions with a location-known statistic and
count lesion-like false blobs in lesion-free reconstructions.

s